\chapter{PENUTUP}

\section{Kesimpulan}
\begin{enumerate}
\item Aplikasi \textit{AI Music Detector} berhasil dibangun sebagai studio web lokal yang menganalisis rekaman musik tanpa mengirim audio ke layanan luar. Setiap vonis disertai skor per segmen 5 detik dan spektrogram yang dilihat model.
\item \textit{Shortcut learning} terhadap codec dicegah dengan augmentasi codec untuk kedua kelas, normalisasi loudness, dan pembatasan spektrogram sampai 7 kHz. Studi ablasi menunjukkan pembatasan 7 kHz adalah komponen kunci: tanpanya, vonis masih berbalik pada separuh lagu yang dikompres ulang.
\item Model final mencapai akurasi 86,5\%, \textit{balanced accuracy} 88,9\%, dan ROC AUC 0,966 pada 3.911 lagu test, dengan salah tuduh hanya 3,5\%. Pada kompresi ulang, akurasinya bertahan di 86--88\% dan hanya 5\% vonis yang berbalik.
\item Kelemahan utama ada pada generator yang tidak ada di data latih: udio-30s hanya terdeteksi 35,1\%, meskipun dua versi Suno yang juga tidak dilatih tetap terdeteksi sekitar 99\%.
\end{enumerate}

\section{Saran}
\begin{enumerate}
\item Menambah contoh lagu dari generator terbaru, terutama Udio dan Suno versi baru, lalu melatih ulang model memakai cache tambahan tanpa memproses ulang seluruh data.
\item Menjalankan evaluasi \textit{leave-one-generator-out} untuk setiap generator agar kemampuan generalisasi terukur lebih lengkap.
\item Mengumpulkan umpan balik pengguna pada lagu yang asal-usulnya jelas sebagai uji lapangan, dengan sebagian lagu disisihkan khusus untuk evaluasi.
\item Mengkalibrasi ambang keputusan dan mengulang \textit{black-box testing} memakai model final.
\end{enumerate}
